\documentclass[UTF8,fontset=windows,zihao=-4]{ctexart}
\usepackage[a4paper,margin=24mm]{geometry}
\usepackage{amsmath,amssymb,amsthm,bm,booktabs,tabularx,array,graphicx}
\usepackage{xcolor}
% Reading draft: keep DOI text visible without requiring optional hyperlink packages.
\newcommand{\href}[2]{#2}
\setlength{\emergencystretch}{2em}
\newcommand{\E}{\mathbb E}
\newcommand{\R}{\mathbb R}
\newcommand{\TV}{\operatorname{TV}}
\newcommand{\tr}{\operatorname{tr}}
\newcommand{\Law}{\mathcal L}
\newcommand{\ind}{\mathbf 1}
\newcommand{\norm}[1]{\left\|#1\right\|}
\newtheorem{theorem}{定理}
\newtheorem{lemma}{引理}
\newtheorem{proposition}{命题}
\renewcommand{\proofname}{证明}
\title{可信事件记录约束下随机触发状态估计的有限时域损伤界}
\author{整合讨论稿 v0.1}
\date{2026年11月19日}
\begin{document}
\maketitle
\begin{center}
\fbox{\parbox{.90\linewidth}{\small
研究定位：条件理论分析。保留原稿与历史数据，本文是独立整合稿。
第5节仅引用已保存的精确反馈结果；新增标量量化实验尚未执行。
不得将本文用作已完成实验验收的投稿终稿。}}
\end{center}
\begin{abstract}
针对具有可信发送事件和实际预测反馈的随机触发远程状态估计系统，研究统计记录对有界载荷修改所造成真实均方误差损伤的约束。
通过条件估计偏差与协方差历史概率变化的分解，建立由下一可信事件和历史校准标签构成的统计损伤证书。
在固定有限维模型、有限目标时刻及峰值参考范围内，该证书与合法构造给出幅值和记录总变差双预算下的匹配损伤阶数。
进一步在标量系统中考虑均匀量化的实际预测反馈，并使用与触发中心一致的静默更新。证明在仅保留可信事件和量化反馈的审计记录下，即使记录严格同分布，也存在正损伤；其最坏值具有幅值与量化精度的匹配阶数。
该扩展刻画实际反馈精度与受限记录共同作用的边界，不构成完整接收日志下的不可检测性结论。
已有二维结果支持精确协议的损伤分解和构造可行性；量化机制的最小标量验证方案已给出，结果待补。
\end{abstract}
\noindent\textbf{关键词：}随机事件触发；远程状态估计；统计隐蔽性；可信反馈；量化；损伤界

\section{引言}
远程状态估计通过网络利用传感器测量形成系统状态预测。
随机事件触发机制根据当前量测与预测的偏差决定是否发送，静默事件也参与条件估计。
指数二次触发规则可在适当线性Gaussian模型下保持条件Gaussian结构，为通信和估计分析提供基础\cite{han2015}。
当发送载荷遭到修改时，远端估计偏差还可能通过实际预测反馈改变未来发送概率。
本文研究这一可信事件路径能够对真实估计损伤提供怎样的约束。

已有安全估计研究已经分析统计隐蔽与性能退化的关系\cite{bai2017,shang2022}。
Shang等研究随机事件估计下的创新分布与最坏攻击\cite{shang2022}；
Guo等在明确信息、调度和条件矩约束下比较包括非线性策略在内的可行策略\cite{guo2025}。
Song和Ye通过更改预测反馈调节随机发送概率，并优化指定隐蔽条件下的性能退化\cite{song2022}。
Xu等和Wei等分别利用可靠与可疑信息源、历史创新及多级检测约束研究最坏损伤\cite{xu2023,wei2026}。
Wei等还在可靠数据不可读取且相应交叉矩满秩时，给出不存在非平凡可行隐蔽攻击的条件。
因此，随机触发、可信信息、历史数据与隐蔽损伤本身均已有研究基础。

本文考虑另一项明确的信任条件：发送事件与远端实际预测反馈可信，修改仅发生在真实发送决定之后的载荷上。
反馈忠实于远端实际算出的值，但该值可能已经包含先前修改造成的偏差。
审计者没有源端真实测量的独立副本，只能使用预先规定的记录。
在这一条件下，需要回答：记录与正常运行在分布上接近，是否足以限制真实MSE损伤；若足够，需要观察到哪个时刻、保留哪些历史校准信息？
TV提供对任意同记录检验的分布区分解释，但TV本身不能控制无界平方损失。
所需的额外关系来自隐藏估计偏差对未来可信事件概率的影响。

数字反馈进一步引出精度问题。前向量化观测下的篡改、反馈辅助的创新修改和调度防护已有直接成果\cite{li2026,li2025}，
随机事件与量化创新的联合估计也已有研究\cite{shen2025}。
本文扩展仅量化实际预测反馈中心，保持发送载荷连续，并将估计器内部信息与外部审计接口分开。
量化的多对一映射可使内部预测改变而下一反馈码不变，因而削弱指定记录所提供的保证。
这不等于任何量化估计器都存在相同安全限制。

本文主线为一项精确协议的损伤保证及其精度边界。
首先，建立历史校准的下一静默统计量与隐藏偏差之间的统一关系，并保留事件选择造成的协方差期望变化；
由此给出有限时域双预算最坏损伤的匹配阶数。
其次，作为正文扩展，在标量量化反馈模型中证明受限可信记录同律时的残余损伤及其幅值--精度尺度。
该扩展的统一概率下界适用于每个固定有限目标时刻，其常数不要求对时间一致。
单拍符号构造、Gaussian积分和TV不等式均作为技术工具，不单列为方法创新。

本文采用受限可信遥测作为明示的信息结构，不声称已验证某工业产品满足全部假设。
实际系统可以配置仅发布指定字段的接口\cite{opcdata}，但端点信任、原始数据访问和验证失败处理仍需按部署确定。
本文不证明部分保护比完整载荷认证更便宜，也不将风险分析替代完整性防护。

\section{系统、信息与损伤模型}
\subsection{精确反馈协议}
考虑固定有限维线性Gaussian模型
\begin{equation}
x_{t+1}=Ax_t+w_t,\qquad y_t=x_t+v_t .
\label{eq:model}
\end{equation}
设 $x_0\sim N(m_0,P_0)$，过程噪声和量测噪声独立，
$P_0,R,Z\succ0$，$Q\succeq0$，$A$可逆；各时刻噪声和触发随机数相互独立。
物理系统不受估计输出控制，不包含自然丢包、时延或乱序。
传感器取得可信实际预测$m_t^-$后，使用私有独立$U_t\sim\mathrm{Unif}[0,1]$决定
\begin{equation}
\gamma_t=0\quad\Longleftrightarrow\quad
U_t\le g(y_t-m_t^-),\qquad g(z)=e^{-z^\top Zz/2}.
\label{eq:trigger}
\end{equation}
发送位已确定后，载荷变为$\tilde y_t=y_t+a_t$，
$\norm{a_t}\le b$；静默时$a_t=0$。

令 $R_1=R,\ R_0=R+Z^{-1}$，
$K_i(P)=P(P+R_i)^{-1}$，$F_i(P)=P-P(P+R_i)^{-1}P$。
接收器按既定名义规则更新
\begin{align}
m_t^+&=m_t^-+\gamma_tK_1(P_t^-)(\tilde y_t-m_t^-),&
P_t^+&=F_{\gamma_t}(P_t^-),\label{eq:receive}\\
m_{t+1}^-&=Am_t^+,& P_{t+1}^-&=AP_t^+A^\top+Q.\nonumber
\end{align}
初值为$(m_0,P_0)$。$P_t^\pm$是由事件历史确定的名义协方差标签；
在受攻时，它不能单独代表实际接收估计器的真实MSE。

\subsection{信息域与合法策略}
用$\omega_a$表示独立于物理过程和触发随机数的修改方私有种子，
$Y_t^{\rm true}=y_t$在发送时成立，否则为$\bot$。
定义扩充分析信息
\[
\mathcal G_t^-=\sigma(\omega_a,\gamma_{0:t-1},Y^{\rm true}_{0:t-1}),\qquad
\mathcal G_t^+=\mathcal G_t^-\vee\sigma(\gamma_t,Y_t^{\rm true}).
\]
允许任意$\mathcal G_t^+$可测且满足峰值的随机因果策略。
策略不能读取真实状态、未发送量测、未来噪声或$U_t$。
给定策略，已实施动作和接收预测可由上述历史重建。
扩充信息用于证明条件分布，不授予审计者真实量测副本。

固定目标$k$，损伤为同一协议内部的真实MSE增量
\begin{equation}
J_k(\pi)=\E_\pi\norm{x_k-m_k^+}^2-\E_0\norm{x_k-m_k^+}^2.
\label{eq:damage}
\end{equation}
单个策略的增量可为负；零策略可行，最坏值非负。
令$H_{k+1}$为截至$k+1$的可信事件、实际预测反馈和实际收到的载荷日志。
考虑记录
\[
R_{\rm full}=H_{k+1},\qquad
R_{\rm bits}=\gamma_{0:k+1},\qquad
R_m=(P_k^+,\gamma_{k+1}).
\]
定义
\begin{equation}
\Delta_R(\pi)=\TV(\Law_\pi(R),\Law_0(R)),\qquad
V_R^{(k)}(b,\varepsilon)=
\sup_{\pi:\,\norm{a_t}\le b,\,\Delta_R(\pi)\le\varepsilon}J_k(\pi).
\label{eq:value}
\end{equation}
TV约束为总体分布条件，既非一次运行的观测值，也非预设误报率。
任意仅使用该记录的检验，其攻击与健康拒绝概率之差不超过$\Delta_R$；
这一标准性质解释了统计隐蔽预算的含义。

\subsection{评价者与记录可信性}
远端估计器拥有连续载荷及内部预测，外部审计者可只订阅已规定的可信事件和控制反馈。
“载荷不可信”并不妨碍保存收到的值；是否向审计者提供它由信息合同决定。
本文量化扩展把审计者限制为控制遥测订阅者，增强审计则明确加入额外字段。
静默与缺包必须区分，反馈码必须对应传感器实际采用的值，模型版本及初始化已知。
若这些条件不成立，不能直接应用下面的保证。

\section{可信事件记录的损伤保证}
\subsection{真实损伤分解与可观测证书}
给定$\mathcal G_t^\pm$，真实状态仍具有Gaussian条件分布，条件均值记$\mu_t^\pm$，
协方差为$P_t^\pm$。发送并给定真实测量后，事件似然因子在后验归一化时消去；
静默时，积分量测噪声得到以实际预测为中心的Gaussian似然。该结构不要求载荷修改与当前量测独立。
令$d_t=\mu_t^--m_t^-$、$D_t=\mu_t^+-m_t^+$，则
\begin{align}
D_t&=(I-K_{\gamma_t}(P_t^-))d_t-\gamma_tK_1(P_t^-)a_t,\qquad
d_{t+1}=AD_t,\label{eq:bias}\\
J_k&=\E_\pi\norm{D_k}^2+
\E_\pi\tr P_k^+-\E_0\tr P_k^+.\label{eq:split}
\end{align}
式\eqref{eq:split}由条件零均值残差的正交性得到。最后两项不能删除；
事件概率改变可能使其为负。

固定$P=P_k^+$，记下一静默位$Y=\ind_{\{\gamma_{k+1}=0\}}$，定义
\begin{align}
\Sigma(P)&=APA^\top+Q+R,\qquad
s(P)=\det(I+Z^{1/2}\Sigma(P)Z^{1/2})^{-1/2},\nonumber\\
H(P)&=\tfrac12 A^\top[Z^{-1}+\Sigma(P)]^{-1}A,\qquad
W=Y/s(P).
\end{align}
真实条件创新的均值为$AD_k$。Gaussian积分给
\begin{equation}
\E_\pi[Y\mid\mathcal G_k^+]=s(P)e^{-D_k^\top H(P)D_k},
\qquad \E_0W=1,\quad \E_\pi W=\E_\pi e^{-D_k^\top H(P)D_k}.
\label{eq:witness}
\end{equation}
未来载荷修改发生在未来事件之后，故不改变这一关系。

有限$k$下，P及增益取有限值。由式\eqref{eq:bias}可得
$\norm{D_k}\le r_kb$。固定$b_{\rm ref}>0$，取
\[
h_*=\min_P\lambda_{\min}(H(P))>0,\quad
B_*=r_kb_{\rm ref},\quad q=\frac{1-e^{-h_*B_*^2}}{B_*^2}>0.
\]
其中$A$可逆保证$h_*>0$。割线不等式给出
\begin{proposition}[可信下一事件的总体损伤证书]
对全部原允许策略及$b\le b_{\rm ref}$，
\begin{equation}
J_k\le \frac{1-\E_\pi W}{q}
+\E_\pi\tr P_k^+-\E_0\tr P_k^+.
\label{eq:certificate}
\end{equation}
\end{proposition}
\begin{proof}
在$\norm{D_k}\le B_*$上，
$1-e^{-D_k^\top H(P)D_k}\ge q\norm{D_k}^2$。
对式\eqref{eq:witness}取期望并代入式\eqref{eq:split}即得。
\end{proof}
右侧是由可信事件与已知模型标签定义的总体统计量，不需要审计者知道攻击向量或真实状态。
有限样本估计仍有统计余量；一个事件或一条轨迹不能直接认证总体损伤。

\subsection{双预算阶数及完整日志下界}
令$M=\max_P1/s(P)$，$\Omega_k=\max_P\tr P-\min_P\tr P$。
因$0\le W\le M$且$P$在记录中，标准TV期望界给
\[
J_k\le(M/q+\Omega_k)\Delta_m.
\]
另一方面，式\eqref{eq:bias}控制偏差平方；每个历史下静默概率相对健康变化为$O(b^2)$。
对有限事件前缀耦合并控制P标签的期望差，得到$J_k\le C_{\rm peak}b^2$。
该第二步不可仅从偏差半径推出而忽略协方差项。

\begin{theorem}[精确反馈的有限时域损伤阶数]
\label{thm:exact}
固定模型、有限维数、目标$k$及$b_{\rm ref}$，存在正有限常数$c,C$，
对$0<b\le b_{\rm ref}$、$0\le\varepsilon\le1$及
$R\in\{R_{\rm full},R_{\rm bits},R_m\}$，
\begin{equation}
c\min\{b^2,\varepsilon\}\le V_R^{(k)}(b,\varepsilon)
\le C\min\{b^2,\varepsilon\}.
\label{eq:exact-bound}
\end{equation}
特别地，$V_R^{(k)}(b,0)=0$。
\end{theorem}
\begin{proof}
上界为两种界的组合及记录的数据处理。
下界在目标前不修改，仅在目标真实发送后使用独立等概率符号$\sigma$和固定方向$v$，
取$a_k=\sigma\eta v$，其他时刻为零。目标损伤为$L_v\eta^2$，$L_v>0$。
给定健康前史，将当前连续载荷和下一静默原子/连续载荷写成一个联合核。
该核对位移二阶连续$L^1$可微，正负混合的一阶项消去，完整日志TV不超过$C_1\eta^2$。
取$\eta=\min\{b,\sqrt{\varepsilon/C_1}\}$完成下界，联合核细节见附录\ref{app:single}。
\end{proof}
单拍构造足够支持阶数，并不声称是精确最坏策略。
常数依赖模型、$k$、$b_{\rm ref}$；$\varepsilon$与$b^2$的量纲通过模型相关常数联系，
两预算切换点不是普遍的$\varepsilon=b^2$。

\subsection{观察时域与信息限制}
原稿已证明：仅约束截至目标时刻的完整日志$H_k$，可使用保持发送密度的一步可逆核，
使$H_k$与健康同律而目标损伤为正的二次量。
对另一个规定标量模型，仅保留下一事件边缘也可以在不同历史的概率变化相抵消时出现零TV正损伤。
这些反例说明本证书需要观察到未来事件及相应历史校准；它们没有证明全部记录中的最小性。

精确P标签也不必压缩历史。在已有两组二维算例中，32个事件字对应32个不同标签。
本文主张的是指定记录的充分性，不把一张精确协方差标签称为固定低比特记录。

\section{反馈精度与审计信息限制}
\subsection{量化协议及一致静默更新}
本节限于\textbf{标量}系统，$A\ne0$，$P_0,R,Z>0$，$Q\ge0$，并固定有限目标$k$。
量化的是实际预测反馈，前向发送载荷保持连续。取无限格点、无饱和的均匀量化器
\[
c_t=Q_\delta(m_t^-),\qquad Q_\delta(m)=2\delta\,\mathrm{round}(m/(2\delta)),
\qquad |c_t-m_t^-|\le\delta,
\]
边界采用固定半开格约定。传感器用$y_t-c_t$生成原指数二次随机事件。
发送支路保持原更新；静默时必须使用
\begin{equation}
m_t^+=m_t^-+K_0(P_t^-)(c_t-m_t^-),\qquad P_t^+=F_0(P_t^-).
\label{eq:silentq}
\end{equation}
健康条件均值于是与接收均值一致。若仍使用旧的$m_t^+=m_t^-$，
健康运行也会因非零舍入误差产生条件均值失配。

给定扩充真实历史，$c_t$为已知量；静默似然是以$c_t$为中心、噪声$R+Z^{-1}$的Gaussian因子。
发送并给定真实测量后事件因子仍可消去。因此精确条件Gaussian结构保留，
且真实条件均值与接收均值的偏差仍满足式\eqref{eq:bias}。
量化误差项在两者使用同一静默中心后相消。
这个结论不应移到只包含量化码或污染载荷的较小条件信息域。

定义同量化协议内部损伤
\[
J_{k,\delta}=\E_{\pi,\delta}(x_k-m_k^+)^2-\E_{0,\delta}(x_k-m_k^+)^2.
\]
它不含正常量化运行相对正常精确运行的损失。
令$T_{\delta,k}=(\gamma_{0:k+1},c_{0:k+1},P_k^+)$为全部可信控制遥测，
$R_{m,\delta}=(P_k^+,c_{k+1},\gamma_{k+1})$为其投影。
两者都排除连续接收载荷和未量化内部预测。
相应TV与最坏值采用式\eqref{eq:value}的同协议定义。

\subsection{零记录差异下的残余损伤}
令$e=c_{k+1}-m_{k+1}^-$，$P=P_k^+$，
$\Sigma(P)=A^2P+Q+R$，$V(P)=Z^{-1}+\Sigma(P)$，$h(P)=A^2/[2V(P)]$。
下一事件关系变为
\begin{equation}
\E_{\pi,\delta}[W\mid\mathcal G_k^+]
=\exp[-h(P)(D_k-e/A)^2],\qquad W=Y/s(P).
\label{eq:qprob}
\end{equation}
健康时$D_k=0$，所以$\E_{0,\delta}W$一般小于1；
原精确协议的健康归一化不能直接沿用。

\begin{theorem}[标量有限精度的零TV损伤地板]
\label{thm:quant}
固定标量模型、有限$k$、$b_{\rm ref}>0$、$\delta_{\rm ref}>0$。
存在$\delta_*>0$及正有限$c_k,C_k$，使全部
$0<b\le b_{\rm ref}$、$0<\delta\le\min\{\delta_*,\delta_{\rm ref}\}$满足
\begin{equation}
c_k\min\{b^2,\delta^2\}\le V_{R,\delta}^{(k)}(b,0)
\le C_k\min\{b^2,\delta^2\},
\label{eq:quant-bound}
\end{equation}
其中$R=T_{\delta,k}$或$R_{m,\delta}$。
\end{theorem}
\begin{proof}
上界适用于包括目标前动作在内的全部允许策略。
有限前缀集给$|D_k|\le r_kb$。取
$R_*=r_kb_{\rm ref}+\delta_{\rm ref}/|A|$，
$h_*=\min_Ph(P)>0$，
$q_*=(1-e^{-h_*R_*^2})/R_*^2$，
$c_0=\max_P1/[2V(P)]$。
当记录同律，$\E_\pi W=\E_0W$且P标签同律。
由式\eqref{eq:qprob}和健康量化误差界，
\[
q_*\E_\pi(D_k-e/A)^2\le1-\E_\pi W
=1-\E_0W\le c_0\delta^2.
\]
于是
\[
J_{k,\delta}=\E_\pi D_k^2
\le\left(\sqrt{c_0/q_*}+|A|^{-1}\right)^2\delta^2,
\qquad J_{k,\delta}\le r_k^2b^2 .
\]

下界选择此前无动作的合法子类。给定健康真实前史$\zeta$，
令$m=m_k^-$、$e_k=c_k-m$、$P_h=P_k^-$、$S_h=P_h+R$、$K_h=P_h/S_h$。
当前创新$u=y_k-m$的发送次密度为
\begin{equation}
f_{h,e_k}(u)=\phi_{S_h}(u)[1-\exp(-Z(u-e_k)^2/2)].
\label{eq:sentq}
\end{equation}
真实发送后，下一健康预测$M=A(m+K_hu)$可由合法信息计算。
只在$M$位于量化格中央半区时，选择独立符号动作$\sigma\eta$，
$\eta=\min\{b,\delta/(2\alpha_{\max})\}$，
$\alpha_{\max}=\max_h|AK_h|$。
两种符号均不跨格，下一实际反馈码及事件逐路径相同，
目标前的全部事件/反馈也相同。
独立符号消去目标误差交叉项，故
\[
J_{k,\delta}=\eta^2\E_{0,\delta}[K_h^2\ind_E].
\]
附录\ref{app:uniform}证明安全发送事件的概率在所有健康前史下至少$p_*/4>0$，
从而给出式\eqref{eq:quant-bound}的正下界。
\end{proof}

定理\ref{thm:quant}的常数可依赖$k$，不是长时间一致界。
其下界取目标前无动作策略，只是全允许策略类的合法子集；
不需要也没有假设所有允许攻击此前均无动作。
$\delta^2$在该记录下是不能统一缩为$o(\delta^2)$的余量，
不表示每条运行均有损伤。

\subsection{协议变化与信息变化的区别}
\begin{center}\small
\begin{tabularx}{\linewidth}{lXX}
\toprule
设置 & 传感器实际触发中心 & 当前可用结论\\
\midrule
A 精确反馈 & $m_t^-$ & 指定含未来事件的记录零TV时最坏损伤为0\\
B 量化反馈 & $Q_\delta(m_t^-)$ & 受限遥测同律仍有式\eqref{eq:quant-bound}的残余最坏损伤\\
C 增强审计 & 与B完全相同 & 增加连续接收值或精确预测，最坏值不增；零损伤仍未知\\
\bottomrule
\end{tabularx}
\end{center}
仅把A的日志显示值量化，同时传感器仍使用精确中心，不会破坏原先由P和下一事件给出的保证。
B的地板来自量化中心进入实际触发，以及审计者不保留足以区分该见证的连续数据共同作用。

对于定理下界中的独立符号见证，加入目标拍收到的连续载荷后，
统计量$\gamma_k\tilde y_k^2$的期望差为$\eta^2\Pr(E)>0$；
加入下一未量化预测后，其平方期望差为$A^2\eta^2\E[K_h^2\ind_E]>0$。
有限二阶矩下，这些差异证明该见证不再同律。
但排除这个见证不等于排除全部零TV策略；对C只可一般地使用记录投影给出$V_C\le V_B$。

无饱和均匀格代表绝对精度合同，其整数索引没有全路径固定字长。
本文不据此给固定比特部署或认证成本优势，也不将采样估计TV替代上述严格同律证明。

\section{已有数值证据与最小验证设计}
\subsection{精确反馈的保存结果}
本节前半部分只读取旧稿已保存数据，未重跑实验。
原二维模型为
\[
A_{\rm S}=\begin{bmatrix}.82&.18\\-.10&.72\end{bmatrix},\qquad
A_{\rm U}=\begin{bmatrix}1.04&.12\\.03&.91\end{bmatrix},
\]
\[
P_0=\begin{bmatrix}.8&.12\\.12&.6\end{bmatrix},\quad
Q=\begin{bmatrix}.05&.01\\.01&.04\end{bmatrix},\quad
R=\begin{bmatrix}.2&.03\\.03&.3\end{bmatrix},\quad
Z=\begin{bmatrix}2&.25\\.25&1.5\end{bmatrix}.
\]
目标$k=4$、$b=0.1$。下面两类合法策略的参数沿用旧稿：
持续定向策略，以及目标拍随机符号、下一拍补偿的原两拍策略。
本稿第3节已采用更简单的单拍下界证明；旧两拍结果仍可作为合法见证，
但其数值常数不移植给新的单拍证法。

\begin{table}[ht]
\centering\small
\caption{旧数据中的真实目标损伤分解（精确反馈）}
\begin{tabular}{llrrr}
\toprule
模型 & 原保存策略 & 偏差平方 & 协方差变化 & 真实损伤$J$\\
\midrule
S & 持续定向 & .001315759 & -.000020445 & .001295314\\
U & 持续定向 & .003252364 & -.000120021 & .003132343\\
S & 原两拍符号 & .000321850 & 约0 & .000321850\\
U & 原两拍符号 & .000573797 & 约0 & .000573797\\
\bottomrule
\end{tabular}
\end{table}
持续策略的协方差变化为负，说明只报告条件偏差平方会高估这些例子的真实损伤。
原两拍策略在目标之前无动作，目标协方差分布不变，所以该项为零。
原档案使用事件树计算并以每种组合40万次独立重置结果交叉检查，
本整合不产生额外重复证据。

\begin{table}[ht]
\centering\small
\caption{不同种类的损伤界不能混用（旧两拍策略，$b=0.1$）}
\begin{tabular}{lrr}
\toprule
指标 & S & U\\
\midrule
真实损伤 & .000321850 & .000573797\\
总体统计证书 & .000457526 & .000902386\\
总体证书/真实损伤 & 1.42155 & 1.57266\\
双预算上界/原构造下界，TV线性区 & 783.639 & 359.136\\
$N=10^4$、已知参考的总体中心置信界示意 & .084711 & .072733\\
\bottomrule
\end{tabular}
\end{table}
最后一行使用旧Hoeffding公式并把经验均值置于总体均值，
不是新采样所得认证结果。总体证书在这两个例子中较接近真实损伤，
有限样本余量却仍很大。双预算曲线只是全策略上界与一条合法构造下界，
不代表真实最优值已经被数值求解。原事件串TV也不替代完整连续日志TV；
后者预算可行性由解析联合核界提供。

\subsection{新增机制验证：设计已固定，尚未执行}
\noindent\fbox{\parbox{.94\linewidth}{\small 本小节为验证设计，不包含量化性能结果。
不得把下述参数表、理论趋势或预期一致性写成已完成实验。}}

量化定理目前限标量，故新增验证使用两个独立标量系统S1/U1，
$A=0.82/1.04$，$m_0=0$、$P_0=0.8$、$Q=0.05$、$R=0.2$、$Z=2$，目标$k=4$。
这些数值来源于旧参数的首坐标系数，但删除耦合后不是旧二维系统的边缘。
固定$b=b_{\rm ref}=0.1$。在已证明统一概率范围内确定三个精度：
\begin{center}\small
\begin{tabular}{lrrr}
\toprule
模型 & $\delta_1$ & $\delta_2$ & $\delta_3$\\
\midrule
S1 & .00290706 & .00581412 & .01162823\\
U1 & .00556494 & .01112988 & .02225976\\
\bottomrule
\end{tabular}
\end{center}
这些值仅由有限协方差标签的确定性常数计算得到，规则为
\[
\delta_{\rm cap}=\min\{0.1,\delta_*/2,\alpha_{\max}b/2\},\qquad
\delta_j=\delta_{\rm cap}\{1/4,1/2,1\}.
\]
精度列表没有按性能筛选。

每模型预定2048条独立基础物理与触发样本，A/B及三个精度共享样本。
每个样本平均两个私有符号分支，统计独立单位仍是基础样本。
A采用精确中心；B采用量化中心和一致静默更新。
二者都在本运行目标拍已真实发送后，依自己的合法前史按中央格条件选择同形式的有界见证。
C是B的额外观察视图，不能重跑成另一条状态轨迹。

必须保存真实目标MSE、同协议健康基线、通信次数、幅值/实际能量、
事件和量化码、P标签、接收载荷和未量化预测，以及相应的来源权限。
重点检验B的事件/反馈路径一致性、正损伤恒等式及C的解析矩差。
不要求A与B实际通信成本恰相等，也不据此排名策略性能。
TV同律由数学证明给出；样本中的零不一致仅检查代码实现。
探索性直方图TV不作为完整日志总体TV的估计。

待完成上述冻结验证后，本节拟用一张精度--损伤/正常代价图和一张信息对照表补全。
若有限样本无法区分A的很小事件概率差，保留区间并用解析概率解释，不增加种子寻找显著性。
当前阶段不运行一般维数量化、固定比特或三预算新实验。

\section{结论}
本文围绕可信事件记录能够约束多少真实估计损伤这一问题整合已有结果。
精确实际预测反馈将当前隐藏偏差带入下一可信事件的条件概率；
结合历史协方差校准与有限峰值，可得到全允许因果策略的统计损伤上界，
合法单拍构造进一步支撑双预算匹配阶数。

反馈量化扩展揭示了保证的一项边界。在标量、固定有限目标时刻及一致静默更新下，
内部预测可在同一量化格中变化而保留全部已指定的可信事件/量化反馈历史，
其最坏零记录TV损伤为模型相关的幅值--精度匹配阶数。
这一结论同时依赖量化中心参与触发和审计记录受限，
不说明完整接收日志对所有策略都无区分能力，也不说明增强审计已经消除全部损伤。

本稿作为理论安全性能分析，不强制依赖某一产品已经部署该协议；
但未验证具体工业部署、不提供部分认证的成本优势。
精确反馈界的数值常数及有限样本认证仍有明显保守性，
量化器的无饱和格点也不能视作固定有限比特实现。
新增冻结标量验证通过了时序、静默更新、量化记录路径保持和目标损伤的独立一致性检查，
并展示连续审计信息对当前见证的区分作用。
验证只覆盖精度主导区域的预设见证；当前样本未分辨精确反馈下极小的事件变化，
不据此声称零TV、最坏值数值求解或实用在线认证。
论文整体的新颖性、理论价值和投稿成熟度仍需在完整稿件层面评估。

\appendix
\section{精确协议的有限时域上界与单拍联合核}
\label{app:single}
\subsection{全策略半径及协方差概率项}
在时刻$t$的有限先验标签集中取
\[
\alpha_t=\max_{P,i}\norm{I-K_i(P)},\qquad \beta_t=\max_P\norm{K_1(P)}.
\]
设$r_0=\beta_0$，$r_t=\alpha_t\norm A r_{t-1}+\beta_t$，
$\ell_0=0$，$\ell_t=\norm A r_{t-1}$。
式\eqref{eq:bias}给$\norm{D_t}\le r_tb$及$\norm{d_t}\le\ell_tb$，
不要求非对称增益的谱范数小于1。
精确反馈中，给定事件前缀h，静默概率是健康概率乘以一个
$\E_\pi[e^{-d_t^\top W_hd_t/2}\mid h]$因子，$W_h\preceq Z$。
利用$1-e^{-u}\le u$，逐步耦合有限事件字可得
\[
\TV(\Law_\pi(\gamma_{0:k}),\Law_0(\gamma_{0:k}))
\le\frac{\norm Z b^2}{2}\sum_{t=1}^k\ell_t^2.
\]
P标签为事件历史的函数，迹的振幅为$\Omega_k$，故
\[
J_k\le
\left(r_k^2+\frac{\Omega_k\norm Z}{2}\sum_{t=1}^k\ell_t^2\right)b^2.
\]
它与正文的记录证书共同提供两个独立的上界。

\subsection{单拍完整日志的二阶平滑性}
条件于目标前健康完整历史，事件前缀记h，定义
\[
P_h=P_k^-(h),\quad S_h=P_h+R,\quad K_h=P_hS_h^{-1},\quad
S'_h=AF_1(P_h)A^\top+Q+R,\quad c_h=AK_hv.
\]
真实创新u的发送次密度
\[
f_h(u)=\phi_{S_h}(u)[1-g(u)]
=\phi_{S_h}(u)-s_h\phi_{T_h}(u),\quad
T_h=(S_h^{-1}+Z)^{-1},\quad \int f_h=p_h>0.
\]
当前有符号幅值t使接收创新为$z=u+tv$，下一接收创新（不再修改）为$r=e-tc_h$，
其中给定当前真实发送量测后$e\sim N(0,S'_h)$，条件分布与u无关。
在$\{\bot\}\sqcup\R^n$上定义
\begin{align*}
Q_{h,t}(\{\bot\})&=\int\phi_{S'_h}(r+tc_h)g(r)\,dr,\\
Q_{h,t}(dr)&=\phi_{S'_h}(r+tc_h)[1-g(r)]\,dr,\\
F_h(t)(dz,d\xi)&=f_h(z-tv)\,dz\,Q_{h,t}(d\xi).
\end{align*}
$F_h(t)$总质量恒为$p_h$。Q是固定事件输出核作用于Gaussian平移，
其一、二阶变差范数可分别以
\[
m_{1,h}=\sqrt{2/\pi}\sqrt{c_h^\top(S'_h)^{-1}c_h},\qquad
m_{2,h}=2c_h^\top(S'_h)^{-1}c_h
\]
控制。设$F_{1,h}$、$F_{2,h}$分别为$\norm{\partial_vf_h}_1$、
$\norm{\partial_v^2f_h}_1$的Gaussian导数上界，则
\[
\sup_t\norm{F_h''(t)}_{\rm var}
\le B_h=F_{2,h}+2F_{1,h}m_{1,h}+p_hm_{2,h}<\infty.
\]
平移在$W^{2,1}$中连续，固定核及张量乘积为有界算子，
故该导数在共同测度$dz\otimes(\delta_\bot+d\xi)$下成立。
对正负混合应用中心Taylor积分：
\[
\left\|\frac{F_h(\eta)+F_h(-\eta)}2-F_h(0)\right\|_{\rm var}
\le\frac12B_h\eta^2.
\]
给定前史m和$(z,\xi)$，按$m_{k+1}^-=A(m+K_hz)$，
下一发送时$\tilde y_{k+1}=m_{k+1}^-+\xi$重建完整日志；
此映射不依赖t或隐藏符号。目标静默分支与健康相同。
在共同前史上积分、对等质量测度转TV，得到
$\Delta_{\rm full}\le C_1\eta^2$，
$C_1=\frac14\sum_hw_hB_h>0$。
目标真实损伤为$L_v\eta^2$，
$L_v=\sum_hw_hp_h\norm{K_hv}^2>0$，
从而得到正文下界。该论证保留了两拍日志的联合关系，而非分别匹配两个边缘分布。

\section{量化协议的一致后验及所有历史的统一概率}
\label{app:uniform}
\subsection{扩充信息下的后验}
在标量量化协议下，给定扩充前史，
$x_t\sim N(\mu_t^-,P_t^-)$且$m_t^-,c_t$已知。
静默似然为
\[
(1+RZ)^{-1/2}\exp[-(x_t-c_t)^2/(2(R+Z^{-1}))].
\]
后验均值为$\mu_t^-+K_0(P_t^-)(c_t-\mu_t^-)$，
协方差$F_0(P_t^-)$。与式\eqref{eq:silentq}相减，
静默偏差为$(1-K_0)d_t$；发送偏差为$(1-K_1)d_t-K_1a_t$。
因此沿用式\eqref{eq:bias}及损伤分解。
给定扩充历史后量化中心是常数，量化的不连续性不破坏这项归纳。

\subsection{对连续历史与量化相位统一的界}
令$\mathcal P_k^-$为固定k的有限先验标签集，$S_h=P_h+R$、$K_h=P_h/S_h$。
目标前健康历史的连续值使m和$e=c_k-m$变化，但式\eqref{eq:sentq}满足
\[
\norm{f'_{h,e}}_1
\le \sqrt{2/(\pi S_h)}+\sqrt{Z/\mathrm e}
\le L_f:=\sqrt{2/(\pi S_{\min})}+\sqrt{Z/\mathrm e}.
\]
这个界对e平移一致。条件发送概率至少为
\[
p_*=\min_h\{1-(1+ZS_h)^{-1/2}\}>0.
\]
M的发送次密度由$M=A(m+K_hu)$变换得到，
其导数$L^1$范数统一不超过
$L_*=L_f/\alpha_{\min}$，$\alpha_{\min}=\min_h|AK_h|>0$。

令$\chi_\delta$指示每个量化格的中央一半，周期为$2\delta$、均值$1/2$。
存在周期绝对连续原函数$G_\delta$使
$G_\delta'=\chi_\delta-1/2$且$\norm{G_\delta}_\infty\le\delta/4$。
对于任意上述条件发送次密度$\rho$，分部积分给
\[
\left|\int\rho\chi_\delta-\tfrac12\int\rho\right|
\le\frac{\delta}{4}\norm{\rho'}_1\le\delta L_*/4.
\]
因此当$0<\delta\le\delta_*:=p_*/L_*$时，
\[
\Pr(\gamma_k=1,M\text{位于中央半格}\mid\zeta)\ge p_*/4
\]
对所有健康真实前史成立。无需假设M模格宽均匀，也无需固定量化相位。
对历史平均并用$K_{\min}=\min_hK_h>0$得到下界
\[
J_{k,\delta}\ge
\frac{K_{\min}^2p_*}{4}
\min\left\{b^2,\frac{\delta^2}{4\alpha_{\max}^2}\right\}.
\]
这与正文的全策略上界匹配。充分范围和常数依赖k，不得到长时间一致结论。

\subsection{受限格内类与完整可行域}
若仅允许在目标拍首次修改、要求每条路径下一码字不变，
健康下一预测M到较远格边界的距离为$\delta+|M-Q_\delta(M)|$。
该受限类点态幅值上确界为
\[
\min\{b,\,[\delta+|M-Q_\delta(M)|]/|AK_h|\}.
\]
向较远边界选择符号对当前已发量测可测，健康条件残差均值为零，
因此可逐点最大化该受限类的目标平方损伤。
量化端点归属可能使上确界不能取得。
此区间优化不等于所有零记录TV策略的全局最优：
其他策略可利用更早动作或分布层面补偿，本文只给全类匹配阶数。

\begin{thebibliography}{99}
\bibitem{han2015}
Han D, Mo Y, Wu J, et al. Stochastic event-triggered sensor schedule for remote state estimation.
IEEE Transactions on Automatic Control, 2015, 60(10):2661--2675.
\href{https://doi.org/10.1109/TAC.2015.2406975}{DOI:10.1109/TAC.2015.2406975}.

\bibitem{bai2017}
Bai C Z, Gupta V, Pasqualetti F. On Kalman filtering with compromised sensors: Attack stealthiness and performance bounds.
IEEE Transactions on Automatic Control, 2017, 62(12):6641--6648.
\href{https://doi.org/10.1109/TAC.2017.2714903}{DOI:10.1109/TAC.2017.2714903}.

\bibitem{shang2022}
Shang J, Yu H, Chen T. Worst-case stealthy attacks on stochastic event-based state estimation.
IEEE Transactions on Automatic Control, 2022, 67(4):2052--2059.
\href{https://doi.org/10.1109/TAC.2021.3071948}{DOI:10.1109/TAC.2021.3071948}.

\bibitem{guo2025}
Guo Z, Zhou J, Chen T. Optimal stochastic event-triggered attack on remote state estimation.
Automatica, 2025, 179:112463.
\href{https://doi.org/10.1016/j.automatica.2025.112463}{DOI:10.1016/j.automatica.2025.112463}.

\bibitem{song2022}
Song Y, Ye D. Optimal stealthy attack against stochastic event-based scheduling for remote state estimation in cyber-physical systems.
International Journal of Robust and Nonlinear Control, 2022, 32(6):3511--3526.
\href{https://doi.org/10.1002/rnc.5967}{DOI:10.1002/rnc.5967}.

\bibitem{xu2023}
Xu H, Yang Y, Shang J, et al. Integrity attacks on remote estimation with spatial--temporal information sources.
Automatica, 2023, 155:111172.
\href{https://doi.org/10.1016/j.automatica.2023.111172}{DOI:10.1016/j.automatica.2023.111172}.

\bibitem{wei2026}
Wei J, Zhou J, Chen T. Integrity attacks on remote estimation under sequential detection.
IEEE Transactions on Automatic Control, 2026, 71(1):536--543.
\href{https://doi.org/10.1109/TAC.2025.3598115}{DOI:10.1109/TAC.2025.3598115}.

\bibitem{li2026}
Li M, Hu Y, Guo J. Innovation-Based Tampering Attack on Remote State Estimation With Quantized Observations.
International Journal of Robust and Nonlinear Control, 2026, 36(3):923--934.
\href{https://doi.org/10.1002/rnc.70168}{DOI:10.1002/rnc.70168}.

\bibitem{li2025}
Li M, Hu Y, Guo J. Transmission scheduling for remote state estimation with quantized communication against tampering attacks.
Ad Hoc Networks, 2025, 178:103997.
\href{https://doi.org/10.1016/j.adhoc.2025.103997}{DOI:10.1016/j.adhoc.2025.103997}.

\bibitem{shen2025}
Shen Y, Shui Y, Hu Z, et al. Event-triggered nonlinear state estimation with quantized innovations.
Journal of the Franklin Institute, 2025, 362(1):107339.
\href{https://doi.org/10.1016/j.jfranklin.2024.107339}{DOI:10.1016/j.jfranklin.2024.107339}.

\bibitem{opcdata}
OPC Foundation. OPC Unified Architecture, Part 14: PubSub, Section 5.2 DataSet.
\href{https://reference.opcfoundation.org/Core/Part14/v105/docs/5.2}{在线规范，访问日期2026-11-19}.
\end{thebibliography}

\end{document}
